\section{Derivative limits and the bias interval}\label{app:modes}

\subsection{Derivative limits}\label{app:derivative-limits}

For $r=0,1$ the bounds and the local limit below are in Trinas and
Valle~\cite[Proposition~2.1 and Theorem~3.1]{tv2026}. We need every
$r$, and a fixed initial history, so we give the proof.

\begin{lemma}[derivative limits]\label{lem:derivative-limits}
Fix $a\in(1/2,1)$ and a finite initial history. Extend $P_n$ by zero
outside $\{0,\ldots,n\}$. Suppose the rescaled
count law converges weakly to a law with smooth density $f$.
For each fixed integer $r\ge0$, linearly interpolate
\[
 \frac{n^{(r+1)a}}{2^{r+1}}\nabla^rP_n(k)
 \quad\hbox{at}\quad (2k-n-r)n^{-a},
 \qquad \nabla P_n(k)=P_n(k)-P_n(k-1).
\]
These interpolants converge locally uniformly to $f^{(r)}$.
\end{lemma}

\begin{proof}
Apply $\nabla^r$ to \eqref{eq:P-rec}. This gives
\[
 \nabla^rP_{n+1}(k)
 =\left(\frac{1+a}{2}-\frac{ak}{n}\right)\nabla^rP_n(k)
 +\left(\frac{1-a}{2}+\frac{a(k-r-1)}n\right)\nabla^rP_n(k-1).
\]
For large $n$ and $1\le k\le n+r$, both coefficients are nonnegative
with sum $1-(r+1)a/n$. For $k=0$ and $k=n+r+1$ one term is zero and
the other has coefficient $(1+a)/2<1-(r+1)a/n$. Hence
$\sup_k|\nabla^rP_{n+1}(k)|\le(1-(r+1)a/n)\sup_k|\nabla^rP_n(k)|$,
and iterating,
\[
 \sup_k|\nabla^rP_n(k)|\le C_{a,r}n^{-(r+1)a}.
\]
Using this at orders $r$ and $r+1$, the interpolants are uniformly
bounded and uniformly Lipschitz. So every subsequence has a further
subsequence that converges locally uniformly. To find the limit, sum
by parts against a smooth function with compact support and use weak
convergence. What comes out is the $r$th distributional derivative of
the limit law. Note that the lattice spacing is $2n^{-a}$, and the
normalization above is exactly its $-(r+1)$ power. Since $f$ is
smooth, the distributional derivative is $f^{(r)}$. All subsequences
have this same limit, and the lemma follows.
\end{proof}


\subsection{Sign patterns}

This part has the sign lemma behind
Theorem~\ref{thm:two-modes}. Notation is as in
Sections~\ref{sec:convex} and~\ref{sec:modes}.
Deleting entries can only remove signs from a list's pattern. An alternating list $p$ satisfies
$p\preceq q$ for every alternating list $q$ with the same first sign and at
least as many entries.

\begin{lemma}[pointwise limits]\label{lem:sp-limit}
Let $v_n$ and $v$ be real functions on $\mathbb R$ with $v_n(y)\to v(y)$ for
every $y$, and let $\Pi$ be an alternating list. If every $v_n$ has values at
most $\Pi$, then so has $v$.
\end{lemma}

\begin{proof}
Let $y_1<\dots<y_j$, and let $I$ be the set of indices $i$ with $v(y_i)\ne0$.
For all large $n$, $v_n(y_i)$ has the sign of $v(y_i)$ for every $i\in I$.
Then $\operatorname{sp}(v(y_1),\dots,v(y_j))$ is the pattern of the list
$(v_n(y_i))_{i\in I}$, which is a subsequence of
$\operatorname{sp}(v_n(y_1),\dots,v_n(y_j))\preceq\Pi$.
\end{proof}

\begin{proposition}[the dominant half]\label{prop:dominant-half}
Under the hypotheses of Proposition~\ref{prop:all-bias-modes}, suppose
$q\ge1/2$. Then $P_n(n-k)\le P_n(k)$ for $2k\ge n$, and the right half has
no strict valley. That is,
\[
 P_n(j)\ge\min\{P_n(i),P_n(k)\}
 \qquad\text{if }i<j<k\text{ and }2i+1\ge n.
\]
Equivalently, its first differences contain no negative value followed
by a positive one. Reflection gives the corresponding left-half result
when $q\le1/2$.
\end{proposition}

\begin{proof}
Write the bias explicitly as $P_{n,q}$ and put
\[
 \Delta_n(k)=P_{n,q}(k)-P_{n,1-q}(k)
       =P_{n,q}(k)-P_{n,q}(n-k).
\]
The forward equation gives
$\Delta_{n+1}(k)=(1-r_n(k))\Delta_n(k)+r_n(k-1)\Delta_n(k-1)$.
At $n=1$, $\Delta_1(1)=2q-1\ge0$. Suppose the assertion holds at
length $n$. At a point strictly above the new center, both parents lie in
the preceding upper half and both coefficients are nonnegative.
At an integer center $\Delta_{n+1}(k)=0$ by reflection. Terms outside the
support vanish. This proves the induction. In particular, the
central first difference of an odd row is nonnegative.

We now keep only the half favored by the start, namely the spatial
indices $2k+1\ge t$ in $Z_t$. At $t=1$ its pattern is contained in $(+,-)$. The central entry is
$2q-1\ge0$, followed by the future right boundary values
$u_s(q)=-\ell_s(1-q)$. Since
$1-q\le1/2$, their pattern is contained in $(+,-)$ by
\eqref{eq:bias-edge-increments}. We now rule out a negative difference followed by a positive one,
also among the future boundary entries. Suppose such a pair first
appeared in the next row. Choose a same-sign parent for each entry.
The preceding proof shows that the first parent comes before the second.
The sole possible parent outside the previous half occurs at the new
central difference of an odd row. That entry is nonnegative, so it
cannot begin the forbidden pair. More explicitly, if a negative entry
at index $i$ of row $t+1$ satisfies $2i+1\ge t+1$, then $2i>t$, so both its parents satisfy $2k+1\ge t$. A positive entry further
right has its parents strictly after these, as before. By induction,
no restricted augmented row contains the pair.
A strict valley would contain such two adjacent differences by
telescoping, which proves the claim.
\end{proof}

\subsection{A start with one step in each direction}

\begin{lemma}[a balanced start]\label{lem:balanced-start}
Fix the first two steps to be $+1,-1$. For $0\le a\le1$, the count
law at each length $n\ge2$ is log-concave, with positive support
$\{1,\ldots,n-1\}$. If $1/2<a<1$, its limiting density
$f_a^{+-}$ is symmetric and log-concave, with
$(f_a^{+-})'(t)<0$ for $t>0$.
Consequently the density conditional on a positive first step satisfies
\[
 h_a'(t)>0\qquad(t<0).
\]
\end{lemma}

\begin{proof}
Starting with one step in each direction removes the boundary obstruction
to the log-concavity induction of~\cite[Proposition~2.2]{glrs2025}.
Write one update as
$g_k=Af_k+Bf_{k-1}$, where
\[
 A=\frac{1+a}{2}-\frac{ak}{n},\qquad
 B=\frac{1-a}{2}+\frac{a(k-1)}n,\qquad d=\frac an.
\]
For $u=f_{k-2}$, $v=f_{k-1}$, $w=f_k$, and $z=f_{k+1}$, expansion gives
\begin{align*}
 g_k^2-g_{k-1}g_{k+1}
 &=(A^2-d^2)(w^2-vz)+(B^2-d^2)(v^2-uw)\\
 &\quad +(A-d)(B-d)(wv-uz)+d^2(w-v)^2.
\end{align*}
For $2\le k\le n-1$ we have $A,B\ge d\ge0$.
The preceding row is log-concave with interval support, so all
three gaps are nonnegative. At the two boundaries one of the
neighboring output masses is zero, and there is nothing to check. We
start from the point mass at one and induct. The recurrence also
shows the masses are strictly positive on the support.

Next we identify the limit. For $t\ge1$, join step $t+1$ to the
earlier step $X_i$ that it copies or reverses, and give this edge the
sign $X_{t+1}X_i$. This builds a random recursive tree on the steps.
The edge signs are independent of the tree and of each other, and each
is $+1$ with probability $(1+a)/2$. Remove the edge from step $2$ to
step $1$. This splits the steps into the part of step $1$ and the part
of step $2$. A new step joins a part with probability equal to the
share of that part among the steps so far. Hence the sizes of the two
parts form a P\'olya urn started from one ball of each color, and the
share of the part of step $1$ converges almost surely to a uniform
variable $U$. Given the sequence of parts that the steps join, each
part is a random recursive tree in its own time, with its own
independent edge signs. So the sum of the steps in the part of step
$1$ is an elephant walk with a positive first step, observed at the
size of that part. The sum in the part of step $2$ is minus an
independent copy of such a walk. Both sizes tend to infinity. By the
almost sure convergence of $S_k/k^a$ given $X_1=1$~\cite[Theorem~3.7]{bercu2018},
\[
 \frac{S_n}{n^a}\longrightarrow
 \mathcal L_a^{+-}:=U^a\mathcal L_a'-(1-U)^a\mathcal L_a''
 \qquad\text{almost surely},
\]
where $\mathcal L_a'$ and $\mathcal L_a''$ are copies of $\mathcal L_a$
and $U$, $\mathcal L_a'$, $\mathcal L_a''$ are independent. With two
positive first steps the same split gives the limit
$\mathcal L_a^{++}=U^a\mathcal L_a'+(1-U)^a\mathcal L_a''$.

These limits have entire densities. Let $\phi_a$ be the characteristic
function of $\mathcal L_a$, and take $1<r<1/a$. Since $ra<1$, we have
$v^{ra}+(1-v)^{ra}\ge1$ for $0\le v\le1$. Hence the bound of
Lemma~\ref{lem:analytic-density} gives
\[
 \bigl|\E e^{it\mathcal L_a^{+\pm}}\bigr|
 \le\int_0^1|\phi_a(tv^a)|\,|\phi_a(t(1-v)^a)|\,dv
 \le4e^{-c|t|^r}.
\]
So, exactly as in the proof of that lemma, Fourier inversion gives
densities $f_a^{+-}$ and $f_a^{++}$ that extend to entire functions.
In particular they are smooth, so Lemma~\ref{lem:derivative-limits}
applies to both starts.

Now interpolate the lattice masses log-linearly. The result is
log-concave. Its integral is within twice the largest atom of the
step interpolation's. The largest atom goes to zero, because the
count laws converge weakly to a law with a continuous density.
Hence the limit is log-concave too. Symmetry follows from the
representation of $\mathcal L_a^{+-}$. Since $h_a>0$, also
$f_a^{+-}>0$. Now $(\log f_a^{+-})'$ is nonincreasing and is zero at zero.
Suppose it were also zero at some $t>0$. Then the density would be
constant on $[0,t]$, which analyticity rules out. Hence
$(f_a^{+-})'<0$ for $t>0$.

Next take two positive initial steps. We claim the first differences
are nonnegative on the negative half-line, and use induction from the
point mass at count two. Each entry on the negative half-line has
both parents in the lower half of the previous row. The one exception is the central
difference of an odd row, and there reflection dominance gives the
sign. The derivative limit of Lemma~\ref{lem:derivative-limits} then
gives $(f_a^{++})'(t)\ge0$ for $t<0$.
Finally
$h_a=((1+a)/2)f_a^{++}+((1-a)/2)f_a^{+-}$.
The second summand has strictly positive derivative there.
\end{proof}

Gu\'erin, Laulin, Raschel, and Simon prove positivity of the conditional
mode for memory sufficiently close to one~\cite[Remark~2.1(a)]{glrs2025}.
We need it for all $a\in(1/2,1)$, in the stronger form of a positive
slope at zero.

\begin{lemma}[the conditional mode is positive]\label{lem:positive-conditional-mode}
For every $1/2<a<1$, we have $h_a'(0)>0$. In particular the unique
conditional mode lies strictly to the right of zero.
\end{lemma}

\begin{proof}
The idea is to bound the reflection difference from below by
something that survives rescaling. Write $f_n(k)=\Prob_+(A_n=k)$ and,
for $z\equiv n\pmod2$,
$d_n(z)=f_n((n+z)/2)-f_n((n-z)/2)$.
With a positive start, the reflection induction of
Proposition~\ref{prop:dominant-half} is strict for $0<z\le n$. Indeed
$d_1(1)=1$, and each noncentral entry of the upper half has a
positive parent with a positive coefficient. The recurrence is
\begin{gather*}
 d_{n+1}(z)=\alpha_n(z)d_n(z+1)+\beta_n(z)d_n(z-1),\\
 \alpha_n(z)=\frac12-\frac{a(z+1)}{2n},\qquad
 \beta_n(z)=\frac12+\frac{a(z-1)}{2n}.
\end{gather*}
Fix $n_0>4a$. Divide by $z$ and multiply by
$\prod_{j=n_0}^{n-1}(1-2a/j)^{-1}$. Call the new values $v_n(z)$.
They satisfy a backward averaging equation, with weights
\[
 q_+=\frac{\alpha_n(z)(z+1)}{z(1-2a/n)},\qquad
 q_-=\frac{\beta_n(z)(z-1)}{z(1-2a/n)}.
\]
In the band we use below, these weights are nonnegative and add up to
one. At $z=1$ the parent $0$ gets weight zero. Let $z_n$ be the
backward chain with these weights. Then
\[
 \E[z_n^2\mid z_{n+1}=z]
 =\frac{n-4a}{n-2a}z^2+\frac{3n-4a}{n-2a}.
\]
With $\gamma_n=\Gamma(n-2a)/\Gamma(n-4a)\sim n^{2a}$, the process
$Y_n=z_n^2/\gamma_n$ has backward drift
$\delta_n=(3n-4a)/\{(n-2a)\gamma_n\}$. Note that
$\sum_{n\ge n_0}\delta_n\to0$ as $n_0\to\infty$.

Fix $R>0$. We start the chain at $z_N=1$ if $N$ is odd and at $z_N=2$
if $N$ is even. We stop it at row $n_0$, or when it leaves the band
$0<z/\sqrt{\gamma_n}\le R$. Since $\sqrt{\gamma_n}\sim n^a$ and $a<1$,
we can take $n_0$ so large that $R\sqrt{\gamma_n}+1\le n$ for all
$n\ge n_0$. Then every state in the band and its two parents lie in
the support, and $\alpha_n(z),\beta_n(z)\ge0$ there. The chain uses the
weights only at band states before it stops, so we need their
nonnegativity only there. The chain moves by $\pm1$, and it can reach
$0$ only from $z=1$, where the parent $0$ has weight zero. So it
cannot leave the band through the bottom, and on leaving the band
$Y>R^2$. Up to the stopping time, the mean of $Y$ is at most
$z_N^2/\gamma_N+\sum_{n\ge n_0}\delta_n$. Taking $n_0$ and then $N$
large, Markov's inequality shows that the chain leaves the band with
probability below one half. Every stopped value of $v$ is positive,
by the strict reflection induction. The band states of row $n_0$ are
the finitely many $z$ with $0<z\le R\sqrt{\gamma_{n_0}}$, and at row
$n_0$ the normalizing product is empty, so $v_{n_0}(z)=d_{n_0}(z)/z$.
Their minimum $m_0>0$ depends only on $n_0$ and $R$, not on $N$. The
chain stops in this row with probability above one half. Since
stopping is bounded, $v_N(z_N)\ge m_0/2$. The
product normalizer is a positive constant times $N^{2a}$ in the
limit. Hence
\[
 \liminf_N N^{2a}\frac{d_N(z_N)}{z_N}>0.
\]
By Lemma~\ref{lem:derivative-limits} for the conditional start,
the first differences converge, and this limit is $4h_a'(0)$. So
$h_a'(0)>0$. Since $h_a$ is unimodal, its mode is to the right of
zero.
\end{proof}

